\documentclass[11pt,a4paper]{article}
\usepackage[T1]{fontenc}
\usepackage[utf8]{inputenc}
\usepackage{mathptmx}
\AtBeginDocument{\let\hbar\relax
  \DeclareMathSymbol{\hbar}{\mathord}{AMSb}{"7E}}
\usepackage{amsmath,amssymb,amsthm}
\usepackage{geometry}
\usepackage{microtype}
\usepackage{hyperref}
\usepackage{booktabs}
\usepackage{enumitem}
\usepackage[numbers]{natbib}
\usepackage{titlesec}

\geometry{top=1.1in,bottom=1.1in,left=1.15in,right=1.15in}

\titleformat{\section}{\normalfont\bfseries\large}{\thesection}{1em}{}
\titleformat{\subsection}{\normalfont\bfseries\normalsize}{\thesubsection}{1em}{}
\titleformat{\subsubsection}{\normalfont\bfseries\itshape\normalsize}
  {\thesubsubsection}{1em}{}
\titlespacing*{\section}{0pt}{14pt}{4pt}
\titlespacing*{\subsection}{0pt}{10pt}{3pt}

\theoremstyle{plain}
\newtheorem{theorem}{Theorem}[section]
\newtheorem{proposition}[theorem]{Proposition}
\newtheorem{corollary}[theorem]{Corollary}
\theoremstyle{definition}
\newtheorem{definition}[theorem]{Definition}
\newtheorem{prediction}[theorem]{Prediction}
\theoremstyle{remark}
\newtheorem{remark}[theorem]{Remark}

\hypersetup{colorlinks=true,linkcolor=black,citecolor=black,urlcolor=black,
  pdftitle={Wave Function Collapse as Irreversible Entropy Production},
  pdfauthor={Joshua F. Sandeman}}

\newcommand{\Pact}{\mathcal{P}_{\mathrm{act}}}
\newcommand{\ket}[1]{|#1\rangle}

% ─────────────────────────────────────────────────────────────────────────────
\begin{document}

\begin{center}
  {\LARGE\bfseries Wave Function Collapse as\\
   Irreversible Entropy Production\par}
  \vspace{6pt}
  {\normalsize Paper~II of the Relational Actualism Series\par}
  \vspace{14pt}
  {\large Joshua F.\ Sandeman$^{1*}$\par}
  \vspace{4pt}
  {\normalsize $^{1}$Independent Researcher, Salem, OR, USA.\par}
  {\normalsize $^{*}$Corresponding author: sansuikyo@gmail.com
   $\cdot$ relationalactualism.org\par}
  \vspace{14pt}
\end{center}

% ── Abstract ─────────────────────────────────────────────────────────────────
\begin{abstract}
The quantum measurement problem---how and when a quantum superposition becomes a
definite classical outcome---has resisted resolution for nearly a century.  We show
that if physical reality is an irreversibly growing directed acyclic graph governed
by the unique self-consistent growth rule established in Paper~I, the measurement
problem does not arise.  An interaction becomes a permanent event in the causal
graph when it produces an irreversible increase in quantum relative entropy
$\Delta S(\rho\|\sigma_0)$ above a specific threshold
$\Delta S^* = 0.601$~nats, computed from the BDG action at the Planck density with
no free parameters.  Below this threshold, interactions are virtual (quantum,
reversible, off-shell).  Above it, they are actual (classical, irreversible,
on-shell).  There is no Heisenberg cut, no observer-dependent collapse, and no
modification of quantum mechanics---only an objective, observer-independent criterion
for when events become real.

The criterion is Lorentz-covariant (proved), causally invariant (Lean~4 verified,
unconditional), and consistent with the Born rule.  We derive three specific,
falsifiable predictions: (i)~a strict null result for gravity-mediated entanglement
experiments (the BMV protocol); (ii)~a parameter-free spin-bath collapse timescale
$t^* = 0.274/g$; and (iii)~a hard ceiling on fault-tolerant quantum array size
$N_{\max} = \eta \times p_{\mathrm{th}}$ (the Kinematic Coherence Bound).  Each
prediction distinguishes this framework from all existing interpretations of quantum
mechanics.
\end{abstract}

%================================================================
\section{Introduction}\label{sec:intro}
%================================================================

The measurement problem is the central conceptual puzzle of quantum mechanics.  The
Schr\"odinger equation describes a world of superpositions: a particle can be in two
places at once, a cat can be simultaneously alive and dead, the moon's position is
indefinite until observed.  Yet every observation, without exception, finds a
definite, particular outcome.  How does the quantum realm of possibilities become the
classical realm of facts?

A century of theoretical effort has produced no consensus.  The Copenhagen
interpretation draws a line---the Heisenberg cut---between the quantum system and
the measuring apparatus, and declines to say what happens at the line.  Many-worlds
posits that every outcome occurs in a branching multiverse, eliminating collapse at
the cost of an enormous ontological commitment.  Spontaneous collapse models (GRW,
CSL) modify the Schr\"odinger equation with stochastic terms, introducing new
parameters.  Pilot-wave theories (de~Broglie--Bohm) add hidden variables.  Each
approach either adds machinery, changes the mathematics, or denies the problem
exists.

This paper takes a different approach, following from the hypothesis established in
Paper~I~\cite{P1}: physical reality is an irreversibly growing directed acyclic
graph~(DAG) whose growth rule is uniquely determined by self-consistency.  If this
hypothesis is correct, the measurement problem \emph{does not arise}, because the
framework contains an objective, observer-independent criterion for when quantum
possibilities become classical facts.  The criterion is not imposed---it is derived
from the unique growth rule.  And it makes specific, falsifiable predictions that
distinguish it from every existing interpretation.

The paper is organised as follows.  Section~\ref{sec:criterion} derives the
actualization criterion from the BDG action.  Section~\ref{sec:measurement} shows how
this dissolves the measurement problem.  Section~\ref{sec:lorentz} proves Lorentz
covariance.  Section~\ref{sec:unruh} resolves the Unruh paradox.
Section~\ref{sec:causal} presents the causal invariance theorem.
Section~\ref{sec:born} discusses Born rule consistency.  Section~\ref{sec:bmv}
derives the BMV null prediction.  Section~\ref{sec:spinbath} derives the spin-bath
collapse timescale.  Section~\ref{sec:kcb} presents the Kinematic Coherence Bound.
Section~\ref{sec:comparison} compares with other approaches.
Section~\ref{sec:conclusions} concludes.

%================================================================
\section{The Actualization Criterion}\label{sec:criterion}
%================================================================

\subsection{On-shell and off-shell in the causal graph}\label{sec:onshell}

In the growing DAG, each vertex represents an actualization event---an irreversible
physical interaction that writes a permanent edge into the causal graph.  Not every
proposed interaction becomes a vertex.  The BDG growth rule (Paper~I,
Section~2~\cite{P1}) assigns an action~$S$ to each proposed vertex based on its
causal neighbourhood.  Vertices with $S > 0$ are accepted; those with $S \leq 0$ are
rejected.

The physical distinction between accepted and rejected vertices corresponds to the
quantum field theory distinction between on-shell and off-shell processes.  An
on-shell process---one satisfying the mass-shell condition
$p^\mu p_\mu = m^2 c^2$---produces a real particle that propagates to a detector.
An off-shell process---a virtual particle exchange---contributes to scattering
amplitudes but does not produce a detectable particle.  In RA, this distinction is
not a matter of mathematical convenience but of ontological status: on-shell
processes are vertices in the graph; off-shell processes are not.

\subsection{The entropic formulation}\label{sec:entropic}

The actualization criterion can be stated in terms of quantum relative entropy.  The
quantum relative entropy $S(\rho\|\sigma_0)$ measures how distinguishable the
state~$\rho$ (produced by the interaction) is from the vacuum state~$\sigma_0$.  An
interaction actualises when:
\begin{equation}
\Delta S(\rho\|\sigma_0) > \Delta S^*,
\label{eq:criterion}
\end{equation}
where $\Delta S^* = 0.601$~nats is the actualization threshold.  This threshold is
computed from the BDG action at the Planck density $\mu = 1$ (Paper~I,
Section~2~\cite{P1}):
\begin{equation}
\Delta S^* = -\ln P_{\mathrm{acc}}(1) = -\ln(0.548) = 0.601~\text{nats} \quad[\text{CV}].
\label{eq:threshold}
\end{equation}

The entropic formulation is primary; the on-shell condition
$p^\mu p_\mu = m^2 c^2$ is its perturbative-regime implementation.  In the
perturbative regime, the two criteria coincide: an on-shell interaction produces a
state sufficiently distinguishable from the vacuum to exceed $\Delta S^*$.  Off-shell
interactions have $\Delta S = 0$ and therefore do not actualise.

\subsection{The threshold is not a free parameter}\label{sec:notfree}

The uniqueness theorem of Paper~I establishes that the BDG growth rule is the only
self-consistent option.  Therefore $\Delta S^* = 0.601$~nats is not one threshold
among many---it is the only threshold compatible with the self-consistency of the
causal graph.  No alternative actualization criterion exists within the framework.

%================================================================
\section{The Measurement Problem Dissolved}\label{sec:measurement}
%================================================================

\subsection{No Heisenberg cut}\label{sec:nocut}

In RA, a quantum state is any state with $\Delta S(\rho\|\sigma_0)$ below the
actualization threshold~$\Delta S^*$.  It has not yet been written into the causal
graph as a permanent event.  An actualization event occurs when, through interaction
with an environment, the relative entropy crosses the threshold.  A new permanent
vertex is written.  Nothing ``collapses''---the wavefunction does not undergo a
sudden, discontinuous change.  What happens is that the causal graph gains a new
vertex, irreversibly.

There is no Heisenberg cut because there is no fundamental distinction between
quantum and classical.  There is only a gradient of actualization density.  The
Schr\"odinger equation is the macroscopic approximation that applies when the
actualization density is low: many proposed events, few actualising.  Classical
mechanics is the macroscopic approximation that applies when the density is high:
most proposed events actualise, behaviour is deterministic.  Both are approximations
to the same underlying discrete process.

\subsection{Schr\"odinger's cat}\label{sec:cat}

Schr\"odinger's cat is not in a superposition of alive and dead states.  A real cat
is a densely self-coupled actualization network---trillions of atoms exchanging
on-shell bosons millions of times per second.  The actualization density is
astronomical.  The cat is classical by constitution, not by observation.  The
radioactive isotope alone is the genuine locus of quantum indeterminacy: a single
system waiting for its first actualization event.  The paradox dissolves because the
thought experiment misidentified the quantum system.  The cat was never quantum;
only the isotope was.

\subsection{The vacuum does not gravitate}\label{sec:vacuum}

A direct consequence of the actualization criterion: off-shell processes have
$\Delta S = 0$ and do not write vertices into the graph.  Since the gravitational
field is sourced only by actualized vertices (via the actualization projector~$\Pact$
applied to the stress-energy tensor), virtual quantum fluctuations do not gravitate.
The cosmological constant is not small---it is identically zero, because the vacuum
contributes nothing to the metric source.  This dissolves the cosmological constant
problem structurally, as a consequence of the same criterion that dissolves the
measurement problem.

%================================================================
\section{Lorentz Covariance}\label{sec:lorentz}
%================================================================

The actualization criterion must be frame-independent: two observers in relative
motion must agree on whether an event has actualised.  This is guaranteed by the
frame-independence of quantum relative entropy.

\begin{theorem}[Frame Independence, LV]\label{thm:frame}
For any unitary $U$ satisfying $U\sigma_0 U^\dagger = \sigma_0$:
\begin{equation}
S(U\rho U^\dagger \| \sigma_0) = S(\rho \| \sigma_0).
\label{eq:frame}
\end{equation}
\end{theorem}

This is proved in Lean~4 (\texttt{frame\_independence} in
\texttt{RA\_AQFT\_Proofs\_v10.lean}).  For Poincar\'e
transformations~$U$, the Minkowski vacuum~$\sigma_0$ is invariant
($U\sigma_0 U^\dagger = \sigma_0$), so the theorem applies.  The
actualization criterion is therefore Lorentz-covariant: if
$\Delta S > \Delta S^*$ in one frame, it exceeds $\Delta S^*$ in all
frames.

This resolves a conceptual difficulty that afflicts many collapse models: the
question of which frame defines the collapse surface.  In RA, no such question
arises.  The criterion is frame-independent by the unitary invariance of relative
entropy.

%================================================================
\section{The Unruh Resolution}\label{sec:unruh}
%================================================================

The Unruh effect~\cite{Unruh1976} poses a challenge to any observer-independent
collapse criterion: a uniformly accelerated (Rindler) observer perceives the
Minkowski vacuum as a thermal bath at the Unruh temperature
$T_U = \hbar a/(2\pi c k_B)$.  If thermal excitations trigger actualization, the
Rindler observer would see events that the inertial observer does not, violating
observer-independence.

\begin{theorem}[Rindler Stationarity, LV]\label{thm:rindler}
Under the modular flow~$\alpha_t$ generated by the Rindler Hamiltonian:
\begin{equation}
\frac{d}{dt}\,S\!\left(\alpha_t(\rho_\beta)\,\|\,\sigma_0\right) = 0.
\label{eq:rindler}
\end{equation}
\end{theorem}

This is proved in Lean~4 (\texttt{rindler\_stationarity} in
\texttt{RA\_AQFT\_Proofs\_v10.lean}).  The KMS thermal state~$\rho_\beta$
associated with the Rindler wedge has identically zero time-derivative of relative
entropy under the modular flow.  The threshold~$\Delta S^*$ is never crossed.  No
actualization occurs.

The physical interpretation: the Rindler thermal bath is a \emph{relabelling} of the
Minkowski vacuum, not a new physical state.  The relative entropy between the
relabelled state and the vacuum is constant---the state is not becoming more
distinguishable from the vacuum over time.  Both observers agree: no physical event
has occurred.  The Unruh ``particles'' are a coordinate artefact, not actualization
events.

%================================================================
\section{Causal Invariance}\label{sec:causal}
%================================================================

The quantum measure for a set of spacelike-separated actualization events must be
independent of the order in which they are processed.  This is the causal invariance
requirement, essential for consistency of the growth dynamics.

\begin{theorem}[Amplitude Locality, LV]\label{thm:amploc}
For the BDG amplitude function, $a(v|C) = a(v|C')$ whenever
$C \cap \mathrm{past}(v) = C' \cap \mathrm{past}(v)$.
\end{theorem}

This is proved in Lean~4 (\texttt{bdg\_amplitude\_locality} in
\texttt{RA\_AmpLocality.lean}) as a theorem of the discrete DAG dynamics, without
any appeal to QFT microcausality.  The proof exploits the transitivity of the
causal order: causal intervals lie entirely in the past, so the BDG action increment
depends only on the causal past of the proposed vertex.

\begin{theorem}[Causal Invariance, LV]\label{thm:causal}
For BDG dynamics, the quantum measure on any set of
spacelike-separated actualization events is independent of the linear
extension:
\begin{equation}
\mu_\sigma(S) = \mu_{\sigma'}(S)
\label{eq:causal}
\end{equation}
for any two linear extensions $\sigma, \sigma'$ of the causal order.
\end{theorem}

This is proved in Lean~4 (\texttt{bdg\_causal\_invariance} in
\texttt{RA\_AmpLocality.lean}), unconditionally, with zero sorry tags.  The proof
combines amplitude locality (Theorem~\ref{thm:amploc}) with the commutativity of
complex multiplication via induction on \texttt{List.Perm}.  No axioms are used.

Causal invariance guarantees that RA's growth dynamics is well-defined regardless of
the observer's foliation of the causal structure---a requirement that many quantum
gravity approaches have difficulty satisfying.

%================================================================
\section{Born Rule Consistency}\label{sec:born}
%================================================================

The Born rule---the probability of a measurement outcome is the squared modulus of
the wavefunction---is consistent with RA's quantum measure but is not derived from
it.  In RA, the Born rule is a constraint on the amplitude function
($a(x) \propto \psi(x)$), not an additional postulate.  The quantum measure
framework (extending Sorkin's classical quantum measure~\cite{Sorkin1994}) assigns a
measure to sets of actualization events; the Born rule specifies which amplitude
function is physically realised.

This is a sharper claim than ``many-worlds derives the Born rule'': RA's quantum
measure is a generalisation of Sorkin's framework, verified to be causally invariant
(independent of the order in which spacelike-separated events are considered)---a
result now machine-checked in Lean~4.  The Born rule consistency is a structural
property of the framework, not a derivation of the rule itself.

%================================================================
\section{The BMV Null Prediction}\label{sec:bmv}
%================================================================

\subsection{The experimental protocol}\label{sec:bmv-protocol}

The Bose--Marletto--Vedral (BMV) protocol~\cite{Bose2017,Marletto2017} proposes to
test whether gravity can mediate entanglement between two masses, each placed in a
spatial superposition.  If the gravitational field is quantum (i.e., capable of
existing in superposition), the two masses should become entangled through their
mutual gravitational interaction.  Detection of this entanglement would constitute
evidence for the quantum nature of gravity.

\subsection{RA's prediction}\label{sec:bmv-prediction}

In RA, the spacetime metric is updated strictly from actualized vertices.  A mass in
quantum spatial superposition---not yet having undergone an actualization event that
commits it to a definite position---has actualization depth $A_{\mathrm{RA}} = 0$
for the position degree of freedom.  It sources zero gravitational field for the
superposed variable:
\begin{equation}
\nabla^2 \Phi = -4\pi G\,\langle\psi_{\mathrm{phys}}|\,\Pact\,
T^{00}\,|\psi_{\mathrm{phys}}\rangle = 0.
\label{eq:bmv}
\end{equation}
No gravitational phase difference accumulates.  No entanglement is generated.

\begin{prediction}[BMV null]\label{pred:bmv}
No gravity-mediated entanglement at any mass scale $m \sim 10^{-14}$~kg,
any superposition duration $\tau \sim 1$~s.  The prediction is strict: a
positive result in any BMV-type experiment would directly falsify RA's
treatment of the gravity-measurement interface.
\end{prediction}

Multiple experimental groups are currently pursuing this test.

%================================================================
\section{The Spin-Bath Collapse Timescale}\label{sec:spinbath}
%================================================================

RA makes a parameter-free prediction for the timescale of wavefunction
actualisation in spin-bath experiments:
\begin{equation}
\boxed{t^* = \frac{1}{g}\sqrt{\frac{\Delta S^*}{2}} = \frac{0.274}{g},}
\label{eq:spinbath}
\end{equation}
where $g$ is the spin-bath coupling constant and $\Delta S^* = 0.601$~nats is the
actualization threshold~[CV].  The derivation: a spin-$\tfrac{1}{2}$ system coupled
to a bath of~$N$ spins with coupling~$g$ accumulates relative entropy at a rate
proportional to~$g^2$.  The threshold~$\Delta S^*$ is crossed at time~$t^*$.

This prediction is specific to RA and distinguishes it from GRW and CSL collapse
models, which predict different functional forms for the collapse timescale (GRW:
exponential in particle number; CSL: dependent on a free parameter~$\lambda$).  The
RA prediction has \emph{no free parameters}: both~$g$ and~$\Delta S^*$ are
determined by the physics of the specific experimental system and the unique BDG
growth rule, respectively.

Superconducting circuit experiments operating at millikelvin temperatures can probe
this regime.  The prediction is that decoherence rates will show a threshold
behaviour at $\Delta S = \Delta S^*$, absent in standard Lindblad master equation
predictions.

%================================================================
\section{The Kinematic Coherence Bound}\label{sec:kcb}
%================================================================

Because each qubit is an additional kinematic site at which actualization events can
occur, RA predicts a structural ceiling on the size of a fault-tolerant quantum array.

\subsection{The bound}\label{sec:kcb-bound}

\begin{prediction}[Kinematic Coherence Bound]\label{pred:kcb}
The maximum number of logical qubits in a fault-tolerant quantum array is
\begin{equation}
N_{\max} = \eta \times p_{\mathrm{th}},
\label{eq:kcb}
\end{equation}
where $\eta$ is the single-qubit quality factor (hardware-dependent) and
$p_{\mathrm{th}}$ is the fault-tolerance threshold (typically $10^{-2}$ to
$10^{-3}$).
\end{prediction}

The derivation: each qubit substrate (electron, photon) has a minimum BDG stability
score of~1~[LV] (\texttt{structural\_fragility} in \texttt{RA\_D1\_Proofs.lean}).
This means each qubit has a nonzero probability per unit time of undergoing a
spurious actualization event---an irreversible quantum-to-classical transition that
constitutes an error.  For~$N$ qubits, the cumulative probability of at least one
spurious actualization scales linearly with~$N$.  When this cumulative probability
exceeds the fault-tolerance budget~$p_{\mathrm{th}}$, the array cannot maintain
coherence regardless of single-qubit fidelity.

\subsection{Distinguishing prediction}\label{sec:kcb-distinction}

Standard decoherence theory predicts continuous (though exponentially difficult)
scaling: in principle, with sufficient error correction overhead, arbitrarily large
quantum computers can be built.  RA predicts a \emph{hard wall} at~$N_{\max}$.  The
distinction is qualitative: RA says the scaling terminates; standard theory says it
merely becomes expensive.  This is testable at the scale of hundreds to thousands of
logical qubits with surface code experiments.

%================================================================
\section{Comparison with Other Approaches}\label{sec:comparison}
%================================================================

\subsection{Copenhagen}\label{sec:copenhagen}

Copenhagen draws the Heisenberg cut and declines to specify the collapse mechanism.
RA replaces the cut with an objective criterion ($\Delta S > \Delta S^*$) that
requires no observer, no measuring apparatus, and no macroscopic/microscopic
distinction.

\subsection{Many-worlds}\label{sec:manyworlds}

Many-worlds eliminates collapse by positing that every branch of the wavefunction is
equally real.  RA eliminates collapse differently: the wavefunction never collapses
because there is no wavefunction in the fundamental ontology---there are only
actualization events in the causal graph.  The ``wavefunction'' is a macroscopic
description that applies when the actualization density is low.  At the fundamental
level, there is only the growing graph with its threshold.

\subsection{GRW / CSL}\label{sec:grw}

Spontaneous collapse models modify the Schr\"odinger equation with stochastic terms
controlled by free parameters ($\lambda$, $r_c$ in GRW; $\lambda$ in CSL).  RA
introduces no free parameters: $\Delta S^*$ is computed from the unique BDG action.
The collapse timescale in RA ($t^* = 0.274/g$) has a different functional form from
GRW/CSL predictions, providing a clean experimental discriminant.

\subsection{de~Broglie--Bohm}\label{sec:bohm}

Pilot-wave theory adds hidden variables (particle positions) guided by the
wavefunction.  RA adds nothing to the ontology---the causal graph is the complete
description.  There are no hidden variables because there is nothing hidden: the
graph is the reality, and its vertices are the events.

\subsection{Summary of distinguishing predictions}\label{sec:summary-pred}

Three predictions distinguish RA from all existing interpretations: (i)~the BMV null
result (not predicted by many-worlds, pilot-wave, or standard Copenhagen); (ii)~the
parameter-free collapse timescale $t^* = 0.274/g$ (not predicted by any
interpretation without free parameters); (iii)~the hard ceiling~$N_{\max}$ on
quantum array size (not predicted by any interpretation that treats decoherence as
continuous scaling).

%================================================================
\section{Conclusions}\label{sec:conclusions}
%================================================================

The measurement problem dissolves within Relational Actualism because the framework
contains an objective, observer-independent criterion for when quantum possibilities
become classical facts: the actualization threshold $\Delta S^* = 0.601$~nats,
derived from the unique self-consistent growth rule of the causal graph
(Paper~I~\cite{P1}).  The criterion is Lorentz-covariant (proved), causally
invariant (Lean~4 verified, unconditional), and consistent with the Born rule.  It
resolves the Unruh paradox, explains why the vacuum does not gravitate, and makes
three specific, falsifiable predictions that distinguish it from every existing
interpretation of quantum mechanics.

The DAG hypothesis is posited.  The actualization criterion is derived.  The
predictions are tests.  If the BMV experiment detects gravity-mediated entanglement,
or if quantum computers scale past the Kinematic Coherence Bound without
encountering a hard wall, the framework is falsified.  If they do not, each
non-observation adds evidence to the hypothesis that physical reality is, at its most
fundamental level, an irreversibly growing causal graph.

Companion papers derive further consequences: coupling constants and particle
spectrum (Paper~I~\cite{P1}), cosmological parameters and nucleation cosmology
(Paper~III~\cite{P3}), and substrate-independent conditions for biological
complexity (Paper~IV~\cite{P4}).  Paper~0~\cite{P0} provides the narrative overview
and complete framework synthesis.

\section*{Acknowledgments}

The formal verification programme uses Lean~4/Mathlib.  AI collaboration (Claude,
Anthropic) contributed to derivation execution, LaTeX compilation, and red-team
review.  The framework is the author's; the execution was collaborative.

% ─────────────────────────────────────────────────────────────────────────────
\begin{thebibliography}{99}

\bibitem{P1}
J.~F.~Sandeman,
``Paper~I: The Physics of a Self-Consistent Causal Graph,''
Paper~I of this series (2026).

\bibitem{P0}
J.~F.~Sandeman,
``Paper~0: A Biography of the Universe,''
Paper~0 of this series (2026).

\bibitem{P3}
J.~F.~Sandeman,
``Paper~III: The Origin, Structure, and Fate of Universes,''
Paper~III of this series (2026).

\bibitem{P4}
J.~F.~Sandeman,
``Paper~IV: The Causal Firewall,''
Paper~IV of this series (2026).

\bibitem{Unruh1976}
W.~G.~Unruh,
``Notes on black-hole evaporation,''
\textit{Phys.\ Rev.\ D} \textbf{14}, 870 (1976).

\bibitem{Sorkin1994}
R.~D.~Sorkin,
``Quantum mechanics as quantum measure theory,''
\textit{Mod.\ Phys.\ Lett.\ A} \textbf{9}, 3119--3128 (1994).

\bibitem{Bose2017}
S.~Bose \textit{et al.},
``Spin entanglement witness for quantum gravity,''
\textit{Phys.\ Rev.\ Lett.}\ \textbf{119}, 240401 (2017).

\bibitem{Marletto2017}
C.~Marletto and V.~Vedral,
``Gravitationally induced entanglement between two massive particles is sufficient
evidence of quantum effects of gravity,''
\textit{Phys.\ Rev.\ Lett.}\ \textbf{119}, 240402 (2017).

\bibitem{BenincasaDowker2010}
D.~M.~T.~Benincasa and F.~Dowker,
``The scalar curvature of a causal set,''
\textit{Phys.\ Rev.\ Lett.}\ \textbf{104}, 181301 (2010).

\bibitem{Yeats2025}
K.~Yeats,
``Combinatorial interpretation of the coefficients of the causal set
d'Alembertian,''
\textit{Class.\ Quantum Grav.}\ \textbf{42}, 145003 (2025);
arXiv:2412.14036.

\bibitem{Rota1964}
G.-C.~Rota,
``On the foundations of combinatorial theory~I: Theory of M\"obius functions,''
\textit{Z.\ Wahrscheinlichkeitstheorie} \textbf{2}, 340--368 (1964).

\end{thebibliography}

\end{document}
